% --- Chapter: Pattern Matching ---
\section{Pattern Matching}

CRISP supports Rust-style pattern matching with the \texttt{match} expression. Patterns can be literal values, variable bindings, or wildcards, with optional \texttt{where} guards.

\subsection{Basic Matching}

\begin{tcolorbox}[colback=codebg, colframe=darkpurple, colbacktitle=darkpurple, title=\textbf{\textcolor{codegold}{Basic Match}}, fonttitle=\bfseries, sharp corners]
\begin{lstlisting}
# Simple value matching
let value = 42;
match value {
    0  => { say "zero"; }
    1  => { say "one"; }
    42 => { say "the answer"; }
    _  => { say "something else"; }
}
# Output: the answer

# Default case with wildcard
let x = 99;
match x {
    1  => { say "one"; }
    2  => { say "two"; }
    _  => { say "unknown: ", x; }
}
# Output: unknown: 99
\end{lstlisting}
\end{tcolorbox}

\subsection{Pattern Types}

CRISP supports the following match patterns:

\begin{table}[H]
\centering
\rowcolors{2}{softviolet!30}{white}
\caption{\textcolor{darkpurple}{Match Pattern Types}}
\vspace{0.3em}
\begin{tabular}{@{}lllp{5cm}@{}}
\toprule
\rowcolor{softvioletdark}
\textbf{\textcolor{white}{Pattern}} & \textbf{\textcolor{white}{Syntax}} & \textbf{\textcolor{white}{Example}} & \textbf{\textcolor{white}{Matches}} \\
\midrule
Integer & \texttt{42} & \texttt{42 => ...} & Exact integer equality \\
String & \texttt{"hello"} & \texttt{"hello" => ...} & Exact string equality \\
Boolean & \texttt{true}, \texttt{false} & \texttt{true => ...} & Exact boolean match \\
Variable & \texttt{name} & \texttt{n => ...} & Any value, binds to name \\
Wildcard & \texttt{\_} & \texttt{\_ => ...} & Any value, no binding \\
\bottomrule
\end{tabular}
\end{table}

\begin{tcolorbox}[colback=codebg, colframe=darkpurple, colbacktitle=darkpurple, title=\textbf{\textcolor{codegold}{Pattern Types}}, fonttitle=\bfseries, sharp corners]
\begin{lstlisting}
# Integer patterns
match 42 {
    0  => { say "zero"; }
    42 => { say "the answer"; }
    _  => { say "other"; }
}

# String patterns
match "hello" {
    "hello" => { say "greeting"; }
    "bye"   => { say "farewell"; }
    _       => { say "unknown"; }
}

# Boolean patterns
match true {
    true  => { say "yes"; }
    false => { say "no"; }
}

# Variable binding — matches anything, binds to name
match readline("Enter: ") {
    n => { say "you entered: ", n; }
}
\end{lstlisting}
\end{tcolorbox}

\subsection{Where Guards}

\texttt{where} guards add conditions to pattern arms. The guard expression is evaluated after the pattern matches, and the arm only fires if the guard is true:

\begin{tcolorbox}[colback=codebg, colframe=darkpurple, colbacktitle=darkpurple, title=\textbf{\textcolor{codegold}{Where Guards}}, fonttitle=\bfseries, sharp corners]
\begin{lstlisting}
# Type-based dispatch
let guess = read("Your guess: ");

match guess {
    n where type(n) == "int" => {
        say "Integer: ", n;
    }
    s where s == "" => {
        say "Empty input.";
    }
    _ => {
        say "Something else: ", guess;
    }
}

# Range-based guards
let score = 85;
match score {
    s where s >= 90 => { say "A"; }
    s where s >= 80 => { say "B"; }
    s where s >= 70 => { say "C"; }
    s where s >= 60 => { say "D"; }
    _               => { say "F"; }
}
# Output: B

# Complex guard expression
let x = 15;
match x {
    n where n > 0 && n % 2 == 0 => { say "positive even"; }
    n where n > 0               => { say "positive odd"; }
    n where n < 0               => { say "negative"; }
    _                           => { say "zero"; }
}
# Output: positive odd
\end{lstlisting}
\end{tcolorbox}

\subsection{Multi-Statement Arm Bodies}

Each arm body can contain multiple statements. The match expression returns the value of the last statement in the matched arm:

\begin{tcolorbox}[colback=codebg, colframe=darkpurple, colbacktitle=darkpurple, title=\textbf{\textcolor{codegold}{Multi-Statement Arms}}, fonttitle=\bfseries, sharp corners]
\begin{lstlisting}
match value {
    1 => {
        say "processing one...";
        say "done!";
    }
    2 => {
        let doubled = 2 * 2;
        say "doubled: ", doubled;
    }
    _ => {
        say "fallback";
    }
}
\end{lstlisting}
\end{tcolorbox}

\subsection{Break and Continue in Match Arms}

\texttt{break} and \texttt{continue} propagate correctly from inside \texttt{match} arms when the \texttt{match} is inside a loop:

\begin{tcolorbox}[colback=codebg, colframe=darkpurple, colbacktitle=darkpurple, title=\textbf{\textcolor{codegold}{Match Inside Loops}}, fonttitle=\bfseries, sharp corners]
\begin{lstlisting}
# break from match inside while
let i = 0;
while true {
    i = i + 1;
    match i {
        n where n >= 5 => { break; }
        _ => { say i; }
    }
}
# Output: 1 2 3 4

# continue from match inside for
for i in 0..5 {
    match i {
        n where n % 2 == 0 => { continue; }
        _ => { say i; }
    }
}
# Output: 1 3
\end{lstlisting}
\end{tcolorbox}

\begin{tcolorbox}[colback=lightlavender, colframe=softvioletdark, title=\textbf{\textcolor{darkpurple}{Design Note: Match Evaluation}}, fonttitle=\bfseries, sharp corners]
  Match arms are evaluated in order from top to bottom. The first arm whose pattern matches and whose guard (if present) evaluates to \texttt{true} wins. Only that arm's body executes. There is no fall-through between arms.
  
For variable-binding patterns with guards, the variable is bound in a temporary scope that includes the guard expression and the arm body. This ensures guard expressions can reference the bound variable without leaking into the outer scope.
\end{tcolorbox}
